% Short-circuit current (opamp-isc): the set-up a person would wire, semantically the deck's.
%   VIsp isp 0 1, XIp isp 0 vp vn osp DUT   +1 V between the inputs: the amplifier saturated high   -> op amp
%   VSCp osp 0 0                            the output to the reference through a 0 V source        -> ammeter I_SC (measured)
%   VIsn, XIn, VSCn                         the same driven negative, sinking                      -> a second copy
%   The swing deck's shorted copies.
\documentclass[tikz, border=10pt]{standalone}
\usepackage{bench}
\begin{document}
\begin{tikzpicture}
\benchtitle{0,7.0}{Short-circuit current}
\def\gnd{1.0}
% the amplifier, non-inverting input up
\node[op amp, noinv input up, font=\small] (A) at (4,4.4) {DUT};
% the supply pins: the rails are what the bench sets
\rails{A.up}{A.down}{$+15$ V}{$-15$ V}
\nodename{anchor=east}{$(A.up)+(-0.08,0.3)$}{vp}
\nodename{anchor=east}{$(A.down)+(-0.08,-0.3)$}{vn}
% the output node
\draw (A.out) -- ++(0.8,0) coordinate (Ao);
\nodename{anchor=south}{$(Ao)+(-0.35,0.05)$}{osp}
% the inverting input grounded
\draw (A.-) -- ++(-0.6,0) coordinate (Apg) -- (Apg |- 0,\gnd) node[ground]{};
\nodename{anchor=south}{$(A.-)+(-0.3,0.03)$}{0}
% the gate: a voltage source from ground
\coordinate (Ai) at (A.+);
\coordinate (Aisrc) at ($(Ai)+(-1.5,0)$);
\draw (Ai) -- (Aisrc);
\draw[set] (Aisrc) to[V, l_={$V_{in}$}] (Aisrc |- 0,\gnd) node[ground, black]{};
\nodename{anchor=south}{$(Ai)+(-0.45,0.02)$}{isp}
% the output shorted to the reference through an ammeter
\coordinate (Asc) at ($(Ao)+(1.5,0)$);
\draw (Ao) -- (Asc);
\ammeter{Asc |- 0,\gnd}{Asc}{$I_{SC}$}
\draw (Asc |- 0,\gnd) node[ground]{};
% the definitions
\definitions{0,-0.1}{%
  $V_{in} = +1$ V (sourcing) and $-1$ V (sinking) between the inputs, one copy each;\quad the output at the reference\\[2pt]
  read $I_{SC}$%ENDDEFS
}
\end{tikzpicture}
\end{document}
